\chapter*{Glosario}
\addcontentsline{toc}{chapter}{Glosario}
\indicetoprule

Este glosario reúne los términos de mercado y decisión empleados de forma recurrente en la memoria.
\begin{description}
  \item[Activo virtual] Elemento digital asociado a una cuenta o a un videojuego. Puede ser funcional, cosmético, coleccionable o consumible.
  \item[Anuncio de venta] Publicación con la que una persona ofrece un activo virtual a un precio determinado.
  \item[Beneficio neto] Resultado económico de una operación después de descontar el precio de compra, las comisiones y los costes aplicables.
  \item[Capital bloqueado] Parte del saldo que permanece comprometida en una compra o que no puede reutilizarse mientras el activo no pueda venderse o transferirse.
  \item[Cartera] Conjunto de activos y saldo disponibles para tomar decisiones de compra, venta o mantenimiento.
  \item[Comisión] Importe o porcentaje que una plataforma descuenta por una operación.
  \item[Conversión de moneda] Transformación de un precio expresado en una divisa a otra para poder compararlo con precios de distintos mercados.
  \item[Exposición] Parte de la cartera comprometida con un activo, una plataforma o una categoría de activos.
  \item[Horizonte temporal] Periodo previsto entre la compra de un activo y su posible venta.
  \item[Liquidez] Facilidad para comprar o vender un activo sin una espera prolongada ni una variación apreciable del precio esperado.
  \item[Mercado cerrado] Entorno en el que los activos solo se utilizan dentro del videojuego y no se intercambian entre personas usuarias.
  \item[Mercado integrado] Mercado administrado por el propio ecosistema del videojuego o de distribución, que controla el inventario y las reglas de intercambio.
  \item[Oferta] Publicación que identifica una variante concreta y fija el precio al que una persona está dispuesta a venderla.
  \item[Operación] Secuencia formada por la compra de un activo y su posible venta posterior.
  \item[Orden de compra] Oferta que indica el precio máximo que una persona está dispuesta a pagar por un activo virtual.
  \item[Plataforma] Servicio digital que administra datos, inventarios, reglas u operaciones relacionadas con activos virtuales.
  \item[Plataforma externa] Plataforma distinta del ecosistema que administra el inventario original y que publica ofertas o facilita operaciones sobre activos transferibles.
  \item[Posición] Activo adquirido que permanece en la cartera hasta su venta o cierre.
  \item[Precio de entrada] Precio publicado para adquirir un activo, expresado en la moneda de comparación.
  \item[Precio de salida] Precio bruto al que se publica o se vende un activo.
  \item[Rentabilidad] Relación entre el beneficio neto de una operación y el capital empleado para realizarla.
  \item[Spread] Diferencia entre el mejor precio de venta y la mejor orden de compra disponibles para un mismo activo en un momento dado.
\end{description}
\indicebottomrule
